% =====================================================================
% ภาคผนวก ก อภิธานศัพท์ — คำเฉพาะที่ผู้อ่านอาจไม่คุ้นเคย เรียงตามตัวอักษร (อังกฤษก่อนไทย)
% =====================================================================
คำศัพท์เฉพาะที่ใช้ในคู่มือและในหน้าจอของระบบมีความหมายดังนี้

\begin{glossarylist}
\term{Compile}{การแปลงไฟล์ \LaTeX{} เป็นไฟล์ PDF}
\term{COMPILE LOG}{บันทึกผลการ Compile แสดงข้อความแจ้งข้อผิดพลาดและบรรทัดที่ผิด}
\term{Docker}{ซอฟต์แวร์สำหรับรันโปรแกรมในคอนเทนเนอร์ที่แยกจากระบบหลัก ใช้รันฐานข้อมูล ส่วนหลัง และ TeX Live}
\term{JWT (token)}{ข้อมูลยืนยันตัวตนที่ระบบออกให้หลังเข้าสู่ระบบ แนบไปกับทุกคำขอ}
\term{LaTeX}{ภาษาสำหรับจัดรูปแบบเอกสาร ใช้เขียนเล่มรายงานโครงงาน}
\term{Nginx}{เว็บเซิร์ฟเวอร์ที่แสดงส่วนหน้าและส่งต่อคำขอไปยังส่วนหลัง}
\term{Snapshot}{สำเนาไฟล์ทั้งหมดของโครงงาน ณ เวลาที่ส่งงาน}
\term{WebSocket}{ช่องทางสื่อสารแบบเรียลไทม์ระหว่างเบราว์เซอร์กับเครื่องแม่ข่าย ใช้ในการแก้ไขเอกสารร่วมกัน}
\term{ข้อเสนอโครงงาน}{เอกสารขออนุมัติหัวข้อโครงงาน ประกอบด้วยชื่อโครงงาน สมาชิก และอาจารย์ที่ปรึกษา}
\term{ฉบับร่าง}{ข้อเสนอที่บันทึกไว้แต่ยังไม่ได้ส่งให้อาจารย์ที่ปรึกษา แก้ไขได้ตลอด}
\term{หัวหน้าโครงงาน}{นักศึกษาผู้สร้างข้อเสนอโครงงาน}
\end{glossarylist}
